\documentclass[11pt]{article}
\usepackage{icmlko}
\begin{document}

\title{규약이 어제 만든 덫에 오늘 내가 걸렸다 --- 봉우리는 열둘이다}
\author{월드모델 실험실}
\date{노트 458 · 2026년 8월 2일}
\maketitle

\begin{abstract}
노트 457 이 깊이 $12$ 를 \emph{집 밖 두 짝을 사고 판을 조금 잃는} 거래로
채택했다. 그러면 바로 물어야 한다 --- \textbf{더 깊으면 더 사나.} 안
물으면 $12$ 는 임의의 자리다.

씨앗 여섯으로 훑으니 \textbf{계속 오르는 것처럼 보였다.}

\begin{center}
\begin{tabular}{lrrr}
\hline
깊이 & 판 & KR 만화 & 비게임 앱 \\
\hline
$12$ (챔피언) & $+0.4839$ & $+0.5561$ & $+0.2666$ \\
$16$ & $+0.4836$ & $\mathbf{+0.5592}$ & $\mathbf{+0.2681}$ \\
$24$ & $+0.4834$ & $+0.5591$ & $+0.2678$ \\
\hline
\end{tabular}
\end{center}

$12 \to 16$ 이 집 밖 두 짝에서 $+0.0031$ 과 $+0.0015$ 인데 \textbf{판은
$-0.0003$ 밖에 안 잃는다}. 노트 457 의 거래보다 \emph{훨씬 좋은} 거래로
보였다 --- 공짜로 전이를 사는 셈이니까.

\textbf{어제 만든 규약이 막았다.} 노트 451 은 \emph{뽑기(또는 씨앗)가
들어가는 비교는 문턱을 잡기 전에 흩어짐을 먼저 재라}고 적었다. 씨앗 열을
\emph{짝지어} 다시 쟀다.

\begin{center}
\begin{tabular}{lrr}
\hline
계기 & 깊이 $12 \to 16$ & 95\% 구간 \\
\hline
판 & $+0.0000$ & $[-0.0012, +0.0013]$ \\
KR 만화 & $+0.0008$ & $[-0.0029, +0.0045]$ \\
비게임 앱 & $-0.0006$ & $[-0.0033, +0.0022]$ \\
\hline
\end{tabular}
\end{center}

\textbf{셋 다 $0$ 을 물고 부호도 갈린다.} 씨앗 여섯이 만든 $+0.0031$ 은
없는 언덕이었다. \textbf{기각 --- 깊이 $12$ 확정.}
\end{abstract}

\begin{figure}[t]
\centering
\includegraphics[width=\textwidth]{figs/f1.pdf}
\caption{\textbf{가} --- 씨앗 여섯으로 그린 곡선. $12$ 를 지나서도
집 밖 두 선이 올라간다. \textbf{나} --- 같은 구간을 씨앗 열로 짝지어
재면 셋 다 $0$ 을 문다.}
\end{figure}

\section{규약이 두 번 일했다}

이 사이클에서 노트 451 의 조항이 두 번 걸렸다.

\textbf{한 번은 나를 고쳤다.} 노트 451 자신이 최근성 이득을 씨앗 넷으로
$+0.0032$ 라 적고 문턱 $+0.003$ 을 넘겨 ``통과''를 찍었는데, 뽑기 흩어짐이
SD $0.0063$ 이라 그 차이가 \textbf{정확히 SE 한 개}였다. 뽑기를 열로
늘려 구간으로 다시 정했다.

\textbf{한 번은 채택을 막았다.} 오늘 깊이 $16$ 이 그것이다. 이번엔
\emph{채택 직전}에 걸렸다는 게 다르다 --- 판 손해가 없어 보였으므로
노트 426 거울조차 필요 없이 넣을 뻔했다.

\textbf{공짜로 보이는 거래를 조심해야 한다.} 노트 457 의 깊이 $12$ 는
판을 확실히 잃었기 때문에 오히려 조심해서 쟀다(씨앗 열 · 짝 12뽑기 ·
네 계기). 깊이 $16$ 은 아무것도 안 잃어 보였기 때문에 하마터면 대충
넣을 뻔했다. \emph{손해가 안 보이면 검사도 안 하게 된다.}

\section{그래서 깊이의 지도가 닫힌다}

\begin{center}
\begin{tabular}{lrl}
\hline
구간 & 집 밖 이득 & 판정 \\
\hline
$6 \to 12$ & $+0.0064 \cdot +0.0061$ (구간 $0$ 밖) & \textbf{채택}(노트 457) \\
$12 \to 16$ & $+0.0008 \cdot -0.0006$ (구간 $0$ 물음) & \textbf{기각} \\
$16 \to 24$ & $-0.0001 \cdot -0.0003$ & 기각 \\
\hline
\end{tabular}
\end{center}

\textbf{봉우리는 열둘이다.} 노트 406 이 깊이 $4 \to 6$ 을 되살리며 적은
문장 --- \emph{봉우리의 자리는 자료 크기의 함수다} --- 이 지금은
$12$ 를 가리킨다.

\section*{남은 것}

\begin{itemize}
\item \textbf{거래 한도가 규약에 없다.} 노트 457 은 판 $-0.0032$ 를
집 밖 $+0.006$ 과 바꿨다. 만약 다음에 판 $-0.02$ 에 집 밖 $+0.04$ 짜리
변경이 오면 넣나? 지금 규약은 \emph{$2\times$짝SD 안}이라는 상한만 있고
\emph{비율}이 없다.
\item \textbf{집 밖 짝 셋째} --- 둘뿐이면 ``둘 다 같은 부호''가 너무
쉽게 찬다.
\item \textbf{lr 은 깊이와 같이 안 봤다} --- lr $0.10$ 은 깊이 $6$ 에서
고른 값이다.
\end{itemize}

\end{document}
